\section{Other proposed ideas}

This section compiles other proposed ideas that aren't directly linked to SVD, but also aren't covered in the exposed literature, and could result in useful results.

\subsection{Weighted minimization}

In previously presented methods, it is assumed that all measurements are equally important. However, this may not be the case. As an example, consider that some of the measurements were taken when the body of water was being disturbed by heavy water currents, affecting the position of the sensors and the perceived sound speed in water. These measurements are useful, but their data is knowingly (\textit{a priori}) less accurate and reliable than that of other measurements, and therefore a weighting step for these measurements can be taken.

Mathematically, this is modeled by adapting \cref{subeq:sec2:iterative_minimization:eq1}. This equation can be expanded, and written as
\begin{equation}
	\bvs{i+1} = \arg \min_{\bvs} \sum_{i=m}^{M} \pts{\bvj[T]{m,i} \bvs - t_{\obs;m}}^2
\end{equation}
in which $\bvj{m,i}$ is the distance traversed in each cell by the $m$-th ray, and $t_{\obs;m}$ is the observed total travel-time for the $m$-th ray. With this formulation, a weighting step can be easily added, resulting in
\begin{equation}
	\bvs{i+1} = \arg \min_{\bvs} \sum_{m=1}^{M} w_m \pts{\bvj[T]{m,i} \bvs - t_{\obs;m}}^2
\end{equation}
where $w_i > 0$ is the weight for the $m$-th ray. Returning to the matricial formulation, we get
\begin{equation}
	\bvs{i+1} = \arg \min_{\bvs} \norm{\bvW \pts{\bvJ{i} \bvs - \bvt{\obs}}}^2
\end{equation}
where $\bvW$ is a diagonal matrix, such that the $\el{\bvW}[m,m] = \Sqrt{w_m}$. Using the SVD and pseudo-inverse approach as established, the solution to this minimization is
\begin{equation}
	\bvG{i} = \pinv*{\bvW \bvJ{i}} \bvW
	\label{eq:sec3:Gi_weighted_minim}
\end{equation}

In the general case this formulation can't be simplified, as the separability of the pseudo-inverse requires $\bvJ{i}$ to have linearly independent rows. This method can be paired with either the regularization or blurring from \cref{subsec:sec2:standard_ttt_approach}, as well as with the null-space exploiting approach as proposed in \cref{subsec:sec3:null-space_exploitation}, by employing the appropriate adaptation of the presented equations.

\subsubsection{Wide matrix weighted minimization}

In the particular case where $N \geq M$ (meaning more cells than than rays) the pseudo-inverse can be distributed, as long as $\rank{\bvJ{i}} = M$. Since $\pinv{\bvW} = \inv{\bvW}$ given it is full-rank, then $\bvG{i} = \pinv{\bvJ{i}}$; meaning the different weighting of measurements isn't relevant for a wide ($N \geq M$) matrix. Another particular case is when $\bvW$ has some zero entries along its main diagonal, leading to the nonexistence of its proper inverse. In this case, $\bvG{i} = \bvJ{i} \pinv{\bvW} \bvW$, where if $\el{\bvW}[m,m] = 0$, then the $m$-th row of $\bvG{i}$ will be zero. That is, since the $m$-th measurement was disregarded ($\el{\bvW}[m,m] = 0$), its weights are $0$ when estimating $\bvs{i+1} = \bvG{i} \bvt{\obs}$.

\subsection{Spatial over-modeling}
\label{subsec:sec4:spatial_over-modeling}

In the literature review portion, it was assumed that $M$ (the number of measurements) is greater than $N$ (the number of discrete cells), to ensure that $\bvJ[T]{i} \bvJ{i}$ was an invertible matrix. However, this requires a large distribution of devices, which isn't necessarily feasible in practice. Therefore, the development/adaptation of theories for an over-modeled (more model variables than measurements) scenario can be useful for cases in which the spatial resolution in the cells is finer than the physical resolution produced by the devices' distribution.

In a scenario where $M < N$,  $\rank{\bvJ[T]{i} \bvJ{i}} = \rank{\bvJ{i}} = R \leq M$, and with $\bvJ[T]{i} \bvJ{i} \in \re^{\sz{N}{N}}$, the product is a rank-deficient matrix, therefore guaranteed to be non-invertible. For the literature-established standard method (\cref{subsec:sec2:standard_ttt_approach}), this forces the use of a regularization matrix and/or of non-inverting zero-gradient solvers.

For the proposed SVD-based inversion, note that at no point it was assumed that $M \geq N$, unlike in the literature review portion. For the proposed structure, $M < N$ ensures that $N-R > 0$, and with this the existence of null-space for $\bvJ{i}$. This ensures the possibility for a minimization on $\bvmu{i+1}$ through the proposed null-space exploiting step. The previously obtained formulation (from \cref{eq:sec2:solution-G_pseudo-inverse,eq:sec2:solution-G_pseudo-inverse_shortform}) is still valid. The fact that $R$ can be equal to $M$ does not change the proposed framework, since at most $\Id{M;R} = \Id$ if $R = M$, and the gradient is still equal to $0$. The null-space exploitation, secondary metric minimization, and SVD truncation processes work exactly the same, but with a guarantee of existence of $\nul{\bvJ{i}}$.

Assuming $R = M < N$, then $\bvU{-R} = []_{\sz{0}{M}}$ is an empty vector, and therefore through \cref{eq:sec2:cost_func_explicit_svd} we have $E(\bvs{i+1}) = 0$. This means that, if all measurements are linearly independent ($R=M$), it is possible to achieve zero error in the measured times. On the contrary scenario, linearly dependent measurements can lead to a non-zero cost function due to value mismatches in the measurements (two measurements that follow the same path but have different times, for example). This is also only valid for the non-truncated SVD scenario, as truncating it will incur in the presence of some residual error.

For the weighted minimization proposed approach, a linear independence in the rows of $\bvJ{i}$ implies that \cref{eq:sec3:Gi_weighted_minim} can be simplified to $\pinv{\bvJ{i}}$. Intuitively, this corresponds to there being more than enough information for the minimization process, and giving different (non-zero) weights to the measurements doesn't affect the result.

Note that over-modeling scenarios, although interesting from the point of having low measurements and high resolution, have to be carefully used. For under-modeled cases, the error of $E(\bvs{i+1})$ comes from the overdetermined system of equation's residues. For over-modeled cases, the error is merely an assessment of the linearly dependence of the measurements, and the mismatch in these. Achieving zero error isn't necessarily a desirable scenario, as it means that for any (linearly independent) solution to the Eikonal equation $\bvJ{i}$, there is a family of slowness fields -- within the null-space of $\bvJ{i}$ -- that fit it perfectly. On the forward problem (used in the FMM/FSM methods for iterating $\bvJ{i}$), this can cause $\bvJ{i+1} = \bvJ{i}$, as there is no error between measured and estimated travel times, without ensuring a convergence to anything similar to the true SSF. With this, inconsistencies between the modeled and true velocity structures may emerge, causing mismatch undetectability by the optimization procedures.

Two possible approaches that avoid this are the use of regularization (with or without blurring) as in \cref{subsubsec:rank-definicient-minimization}, and the truncation of $\bvSi$ from \cref{subsec:sec3:svd_truncation}. Both these techniques would lead to some error in the primary cost function (see \cref{eq:sec3:cost_function_truncatedSVD,eq:sec3:Delta_cost_function_truncatedSVD} for the truncated SVD in particular), leaving space for the iterative Eikonal optimization to find more suitable curves.

\subsection{Perturbation modeling}

We now assume that the slowness $\bvs{i}$ is composed of constant term $\bvs{0}$, and a perturbation term $\bvs[t]{i}$, as
\begin{equation}
	\bvs{i} = \bvs{0} + \bvs[t]{i}
\end{equation}

With this,
\begin{equation}
	\begin{split}
		\bvs{i+1}
		& = \arg \min_{\bvs} \norm{\bvJ{i} \bvs - \bvt{\obs}}^2 \\
		& = \arg \min_{\bvs[t]} \norm{\bvJ{i} \bvs[t] - \pts{\bvt{\obs} - \bvJ{i}\bvs{0}}}^2
	\end{split}
\end{equation}
whose solution is
\begin{equation}
	\bvs{i+1} =  \pinv{\bvJ{i}} \bvt{\obs} - \bvV{2}\bvV[T]{2} \bvs{0}
	\label{eq:sec4:bvz_perturbation_model}
\end{equation}
where $\bvV{2}$ are the last $R$ columns of $\bvV$ corresponding to the null-space of $\bvJ{i}$ (see \cref{subsec:sec3:null-space_exploitation}).

Comparing with \cref{eq:sec2:definition_zi+1_with_mu}, we have that these formulations are very similar, with $\bvD = \Id$ and (supposing that) $\pinv{\bvJ{i}} \bvt{\obs} = \bvs{0}$. Note that the constant term $\bvs{0}$ doesn't need to be constant over the entire body of water, and could be constructed by taking the average over sections of the body of water from previous measurements and estimations.

\subsubsection{Regularization}

This framework can be paired with a regularization term, with it being applied on either the slowness model $\bvs{i+1}$, or on the perturbation model $\bvs[t]{i}$. For the former (regularization on slowness), this leads to
\begin{equation}
	\bvs{i+1} = \arg\min_{\bvs[t]} \norm{\bvJ{i} \bvs[t] - \pts{\bvt{\obs} - \bvJ{i} \bvs{0}}}^2 + \alpha^2\norm{\bvD \pts{\bvs[t] + \bvs{0}}}^2 
\end{equation}
whose solution is given by
\begin{subequations}
	\begin{gather}
		\bvs[t]{i+1} = \inv*{\bvJ[T]{i} \bvJ{i} + \alpha^2 \bvD[T] \bvD} \bvJ[T] \bvt{\obs} - \bvs{0} \\
		\bvs{i+1} = \inv*{\bvJ[T]{i} \bvJ{i} + \alpha^2 \bvD[T] \bvD} \bvJ[T] \bvt{\obs}
	\end{gather}
\end{subequations}
which leads to the same solution as using the traditional regularization in \cref{eq:sec2:solution_regularization-problem_simple}, without the perturbation modeling step.

For the latter (regularization on perturbation), the formulation is
\begin{equation}
	\bvs{i+1} = \arg\min_{\bvs[t]} \norm{\bvJ{i} \bvs[t] - \pts{\bvt{\obs} - \bvJ{i} \bvs{0}}}^2 + \alpha^2\norm{\bvD \pts{\bvs[t]}}^2 
	\label{eq:sec3:minimization_perturbation-with-regularization_formulation}
\end{equation}
which results in
\begin{subequations}
	\begin{gather}
		\bvs[t]{i+1} = \inv*{\bvJ[T]{i} \bvJ{i} + \alpha^2 \bvD[T] \bvD} \pts{\bvJ[T]{i} \bvt{\obs} - \bvJ[T]{i} \bvJ{i} \bvs{0}} \\
		\bvs{i+1} = \inv*{\bvJ[T]{i} \bvJ{i} + \alpha^2 \bvD[T] \bvD} \bvt{\obs} + \pts{\Id{N} - \inv*{\bvJ[T]{i} \bvJ{i} + \alpha^2 \bvD[T] \bvD} \bvJ[T]{i} \bvJ{i}} \bvs{0}
	\end{gather}
\end{subequations}

The first term here is identical to the original regularization problem's solution in \cref{eq:sec2:solution_regularization-problem_simple}, and the second-one gives a correction that penalizes deviating from the constant term: $\alpha = 0$ results in the non-regularized solution, and $\alpha \to \infty$ leads to $\bvs{i+1} = \bvs{0}$. Thus, $\alpha$ can be seen as an interpolation control parameter between the two.

\subsubsection{Null-space exploiting}

Although this formulation based on perturbation upon the average behavior can be useful, its application allied with the proposed null-space exploiting would render the perturbation modeling useless, given that $\bvV{2}\bvV[T]{2} \bvs{0}$ is within the null-space of $\bvJ{i}$, and this term would be dismissed in the minimization process for any secondary given metric.

Mathematically, \cref{eq:sec3:cost_function_mu} with \cref{eq:sec4:bvz_perturbation_model} results in
\begin{equation}
	F(\bvmu{i+1}) = \norm{\pinv{\bvJ{i}} \bvt{\obs} - \bvV{2}\bvV[T]{2} \bvs{0} + \bvV{2} \bvmu{i+1}}^2
\end{equation}
whose minimization leads to
\begin{equation}
	\begin{split}
		\bvmu{i+1}
		& = \inv*{\bvV[T]{2} \bvD[T] \bvD \bvV{2}} \bvV[T]{2} \bvD[T] \bvD \bvV{2} \bvV[T]{2} \bvs{0} - \pinv*{\bvV[T]{2} \bvD[T] \bvD \bvV{2}} \bvV[T]{2} \bvD[T] \bvD \pinv{\bvJ{i}} \bvt{\obs} \\
		& = \bvV[T]{2} \bvs{0} - \inv*{\bvV[T]{2} \bvD[T] \bvD \bvV{2}} \bvV[T]{2} \bvD[T] \bvD \pinv{\bvJ{i}} \bvt{\obs}
	\end{split}
\end{equation}
where the inverse is taken given that both $\bvV{2}$ and $\bvD$ are full-rank matrices, by definition. From this,
\begin{equation}
	\begin{split}
		\bvs{i+1}
		& = \pinv{\bvJ{i}} \bvt{\obs} - \bvV{2}\bvV[T]{2} \bvs{0} + \bvV{2} \bvV[T]{2} \bvs{0} - \bvV{2}\inv*{\bvV[T]{2} \bvD[T] \bvD \bvV{2}} \bvV[T]{2} \bvD[T] \bvD \pinv{\bvJ{i}} \bvt{\obs} \\
		& = \pinv{\bvJ{i}} \bvt{\obs} - \bvV{2}\inv*{\bvV[T]{2} \bvD[T] \bvD \bvV{2}} \bvV[T]{2} \bvD[T] \bvD \pinv{\bvJ{i}} \bvt{\obs}
	\end{split}
\end{equation}
which is identical to \cref{eq:sec2:definition_zi+1_with_mu}, after some algebraic manipulations and appropriate use of the pseudo-inverse's properties. Therefore, the application of a secondary minimization step nullifies the effects of perturbation modeling. Intuitively, by definition the secondary minimization traverses the null-space of $\bvJ{i}$ searching for a minimum in the secondary metric, and the addition of the $\bvV{2}\bvV[T]{2} \bvs{0}$ term would only dislocate the solution by this value, canceling with its addition in \cref{eq:sec4:bvz_perturbation_model}; exactly as seen here.

%\subsection{Device position uncertainty}
%
%In \cite{hellmann_borehole_2023}, a framework that jointly estimates the velocity field and the devices' positions was proposed, where some assumptions on the devices' displacement were taken. 
